\documentclass[UTF8,fontset=none,11pt,a4paper]{ctexart}
\usepackage[left=21mm,right=21mm,top=27mm,bottom=22mm,headheight=26pt,headsep=6mm]{geometry}
\usepackage{fontspec}
\setmainfont{texgyretermes-regular.otf}[BoldFont=texgyretermes-bold.otf,ItalicFont=texgyretermes-italic.otf]
\setCJKmainfont{Songti.ttc}[Path=/System/Library/Fonts/Supplemental/,FontIndex=6,AutoFakeBold=2]
\setCJKsansfont{Songti.ttc}[Path=/System/Library/Fonts/Supplemental/,FontIndex=6,AutoFakeBold=2]
\usepackage[table]{xcolor}
\definecolor{TableHead}{HTML}{EEF0ED}
\definecolor{Muted}{HTML}{666C68}
\usepackage{booktabs,longtable,array,fancyhdr,needspace,xurl}
\usepackage[unicode,hidelinks]{hyperref}
\urlstyle{same}
\linespread{1.32}
\setlength{\parindent}{2em}
\setlength{\parskip}{3pt}
\setlength{\LTpre}{8pt}\setlength{\LTpost}{10pt}
\ctexset{section={format=\fontsize{15}{20}\selectfont\bfseries,beforeskip=18pt,afterskip=8pt},subsection={format=\fontsize{12}{16}\selectfont\bfseries,beforeskip=12pt,afterskip=6pt},subsubsection={format=\normalsize\bfseries,beforeskip=9pt,afterskip=4pt}}
\pagestyle{fancy}\fancyhf{}
\renewcommand{\headrulewidth}{.3pt}
\renewcommand{\footrulewidth}{0pt}
\fancyhead[L]{\parbox[b]{\linewidth}{\fontsize{8}{10}\selectfont\color{Muted}住房贷款：固定月供还是先固定后浮动·浮动利率借据B}}
\fancyfoot[C]{\small\color{Muted}\thepage}
\hypersetup{pdftitle={浮动利率借据B},pdfsubject={仅供模拟使用·住房贷款：固定月供还是先固定后浮动·17-02},pdfauthor={金融决策模拟材料}}
\begin{document}
{\noindent\fontsize{20}{25}\selectfont\bfseries 浮动利率借据B\par}
\vspace{7pt}{\noindent\small\color{Muted}仅供模拟使用·住房贷款：固定月供还是先固定后浮动·17-02\par}\vspace{10pt}
\section*{1. 借款人的付款承诺}
本人已从贷款机构取得216,000美元本金，承诺按贷款机构指示偿还本金及利息。所有本票据项下款项将以美元，采用现金、支票、汇票或贷款机构接受的其他付款方式支付。


贷款机构可以转让本票据。贷款机构，或通过转让取得本票据并有权收取有关付款的人，称为“持票人”。


\section*{2. 利息}
未偿本金将持续计收利息，直至本金全部清偿。本人初始按年利率3.00\%支付利息；利率按本票据第4条调整。


第2及第4条要求的利率，适用于第7(B)项所述违约发生之前及之后。


\section*{3. 付款}
\textbf{(A) 时间及地点。}本人每月支付本金及利息，称为每月付款额。付款自贷款记录列明的首期日期开始，在每月指定日期进行，直至本金、利息及本票据项下其他费用全部支付。


每月付款按预定到期日应用，先用于利息，再用于本金。三十年合同到期日仍有应付款项时，该余额于到期日一次付清。付款寄往持票人指定地点；持票人可以通知其他付款地点。


\textbf{(B) 初始每月付款额。}最初六十个月，每月本金及利息付款约910.66美元；之后金额可以变动。该金额不包括房产税、保险及其他每月费用。


\textbf{(C) 月供变动。}未偿本金及适用利率变动将反映在月供中。持票人按第4条决定新的利率及月供金额。


\section*{4. 浮动利率及每月付款额变动}
\textbf{(A) 调整日。}初始固定利率在第六十一期付款开始时转为浮动利率，此后每十二个月可以调整一次。初始利率转为浮动利率的日期，以及之后每个可以调整利率的日期，称为“调整日”。


\textbf{(B) 指数。}自首个调整日起，利率采用向公众公布的一年期美国国债指数。合同观察日的适用指数，称为“当前指数”。本次合同采用各调整日前四十五日已经公开可得的最近指数；指数为负时，在利率计算中按零处理。


指数停止提供时，按第4(G)项更换。指数的未来数值未在本票据中固定。


\textbf{(C) 计算调整。}调整日前，持票人在当前指数上加2.50个百分点的合同加点。利率按指数及加点的实际合计计算，不按八分之一个百分点取整。第4(D)项的上限及下限仍须应用。


持票人按调整时预计未偿本金、新利率及剩余摊还期数，计算足以在合同到期日清偿本金的大致相等每月付款，作为新的月供。


\textbf{(D) 调整限制。}首次及以后每次调整，相对之前适用利率最多上升或下降2.00个百分点；所有调整后利率均不得低于2.50\%，也不得高于8.00\%。这些限制同时适用，不因指数加点计算结果超出范围而失效。


\textbf{(E) 生效日。}新利率在相应调整日生效；新月供按付款记录中对应新利率的第一期开始支付，并持续至下一次付款额变更。


\textbf{(F) 变动通知。}持票人在生效日前交付或邮寄利率变动通知，列明新的月供、法律要求的资料，以及可以回答通知相关问题的联系人的职务和电话。


\textbf{(G) 替代指数及加点。}指数管理机构永久或无限期停止向公众提供指数，或管理机构及其监管者正式宣布指数不再可靠或具有代表性时，构成替代事件。


发生替代事件时，持票人选择替代指数，必要时选择替代加点。联邦储备委员会、纽约联储或其认可组织已为包括住房浮动利率贷款在内的消费者产品选定或推荐替代指数时，优先采用该指数。


没有上述推荐指数时，持票人应合理且善意地考虑原指数与替代指数的历史表现，选择预计能尽量减少贷款成本变化的替代指数及加点。


替代指数及适用替代加点在替代事件后启用，用于该事件四十五日之后发生的调整日。贷款期间可以再次更换，但须再次发生替代事件。更换后，本文对指数及加点的引用相应指向替代值。


持票人应通知借款人替代指数、适用替代加点，以及法律和监管要求的其他资料。


\section*{5. 借款人的提前还款权利}
本人可以在本金到期前支付本金；仅偿还本金的付款称为提前还款。提前还款时，以书面形式告知持票人。未付清当时已经到期的所有月供时，不得指定一笔付款为提前还款。


本人可以全部或部分提前还款，不支付提前还款收费。持票人可以先从该款项支付对应金额的应计及未付利息，再用余额减少本金。


部分提前还款不改变付款到期日，除非持票人书面同意。提前还款后的首个利率调整日，月供可能降低，但利率提高也可能抵消该降低额。


\section*{6. 贷款费用}
适用法律设有最高贷款费用，且经最终解释，有关已收或将收利息及费用超过限额时，应减少收费至合法限额，并退还已经收取的超额金额。


持票人可以直接退款，也可减少未偿本金。以退款减少本金时，作为部分提前还款处理。


\section*{7. 借款人未按要求付款}
\textbf{(A) 逾期收费。}合同指定宽限期结束时尚未收到完整月供，持票人可以按约定百分比收取逾期费用。本人及时支付，每笔逾期付款只收取一次。


\textbf{(B) 违约。}本人未在每月付款到期日支付完整月供，构成违约。


\textbf{(C) 通知。}持票人可以书面通知：未在指定日期补足逾期金额时，将要求立即支付未偿本金、应付利息及其他本票据费用。该补付日期至少在通知邮寄或以其他方式送达后三十日。


\textbf{(D) 不放弃权利。}某次违约后未立即要求全额清偿，不表示持票人放弃对持续违约或再次违约要求立即清偿的权利。


\textbf{(E) 执行费用。}依法要求立即清偿时，借款人须支付法律允许的执行成本，包括合理律师费及其他成本。


\section*{8. 通知}
\textbf{(A) 给借款人。}适用法律未要求其他方式时，通知以交付或普通邮件送往物业地址，或借款人已通知的新地址。借款人应及时通知实际及邮寄地址变更。


借款人及持票人书面同意且已提供有效电子地址时，可以使用电子邮件或其他电子通信；借款人须及时更新该电子地址。


\textbf{(B) 给持票人。}借款人发出的有关通知，以普通邮件寄往指定收款地址；已经接到地址变更通知的，寄往新地址。


\section*{9. 本票据下各人的义务}
多位签署人中，每个人均个人承担所有承诺，包括支付本票据下全部应付款项。保证人、担保人、背书人，以及承接这些义务的人，同样承担有关责任。


持票人可以向任一个人单独执行，也可以向全部义务人一并执行；任何一个人可能被要求支付全部应付金额。


\section*{10. 权利放弃}
本人及其他义务人放弃提示付款及拒付通知的权利。提示付款是要求持票人正式提出支付到期金额要求的权利；拒付通知是要求持票人通知其他人有关金额未获支付的权利。


\section*{11. 统一担保票据}
本票据采用统一文件结构，某些司法辖区存在有限差异。与本票据同日并经固定／浮动利率附加条款修改的抵押、抵押契据、信托契据或担保契据，合称“担保文件”，保护持票人免于本人未遵守付款承诺可能造成的损失。担保文件还说明可以要求立即清偿全部应付款项的情形和条件。


\textbf{(A) 初始固定利率期间。}在第4条规定的初始固定利率转为浮动利率之前，担保文件第19条适用以下规定。


就本条而言，“物业权益”包括物业的法律权益或受益权益，包括借款人有意在未来将产权转给买受人的契据保证、契据合同、分期销售合同或托管协议项下转移的受益权益。


未经贷款机构事先书面同意，物业或任何物业权益被全部或部分出售或转让，或非自然人借款人的受益权益被转让时，贷款机构可以要求立即全额支付由担保文件保障的全部金额；适用法律禁止时，贷款机构不得行使该权利。


贷款机构行使该权利时，将发出加速清偿通知。自担保文件第16条规定的通知送达之日起，应给予至少三十日清偿全部担保金额。借款人未在期间届满前或届满时付清，贷款机构可以行使担保文件允许的救济，无须另行通知或要求，并可收取寻求救济产生的费用，包括合理律师费、物业检查和估价费，以及保护其物业权益或担保文件权利的其他费用。


\textbf{(B) 浮动利率期间。}初始固定利率依第4条转为浮动利率时，第11(A)项规定停止适用，担保文件第19条改为以下规定。


就本条而言，物业权益仍包括物业的法律或受益权益，以及契据保证、契据合同、分期销售合同或托管协议中旨在未来将产权由借款人转给买受人的受益权益。


未经事先书面同意出售或转让物业、物业权益，或非自然人借款人的受益权益时，贷款机构可以要求立即全额清偿担保金额；适用法律禁止时，不得行使。


借款人向贷款机构提交评价拟受让人所需资料，按向该受让人发放新贷款的标准进行评价，且贷款机构合理判断贷款承接不会损害担保权益、违反担保文件承诺的风险可以接受时，贷款机构也不行使该加速清偿选择权。


在适用法律允许的范围内，贷款机构可以收取合理费用，作为同意承接贷款的条件；也可以要求受让人签署可接受的承接协议，承担本票据及担保文件中的全部承诺。未获贷款机构书面解除之前，原借款人继续承担本票据及担保文件下的义务。


贷款机构要求立即全额支付时，应给予至少三十日加速清偿通知。未在该期限内支付，贷款机构可以行使担保文件允许的救济并收取有关费用，包括合理律师费、物业检查和估价费，以及保护其物业权益或担保文件权利的其他费用。


\textbf{签署人：}［借款人］。本中文文本为参考文件，不供签署。

\end{document}